\chapter{系统误差}
\label{chap:systematics}

\section{概述}
\label{sec:syst-overview}

在第5章信号提取与效率修正的基础上，本章讨论各项系统误差来源。目前已定量传播至结果的分量为采用不同子过程MC样本构造效率图时得到的效率修正变化（第\ref{subsec:syst-model-dependence}节）；积分亮度和分支比的不确定度列于第\ref{sec:syst-normalization}节，其余误差源的方法描述和定性讨论见本章后续各节，尚待完成定量验证与统一传播。因此，本章给出的误差汇总不代表完整的实验系统误差。系统误差的讨论围绕分析链条的各个环节组织，本分析考虑以下类别的系统效应：

\begin{itemize}
  \item 与信号提取程序相关的误差（信号和本底质量模型、固定拟合参数、拟合范围）；
  \item 接受度和效率修正中的误差（MC 统计误差、分箱效应、缪子和带电强子重建与鉴别效率、触发效率、顶点拟合和事例级选择效率）；
  \item 与候选重建和选择相关的误差（多重候选处理、最佳候选选择、候选匹配和迁移效应）；
  \item 样本建模相关的误差（信号 MC 建模、运动学重加权、效率修正的模型依赖性、事例堆积建模）；
  \item 归一化误差（积分亮度、分支比）。
\end{itemize}

已评估的误差源已传播至修正后信号产额和微分分布；其余误差源在后续各节中保留方法框架，供未来定量研究参考。

\section{信号提取相关的系统误差}
\label{sec:syst-signal-extraction}

信号提取的系统误差通过在一组合理的替代方案中改变拟合模型，记录提取的原始信号产额 $N_{sss}$ 的变化来评估。以下各小节列出本分析已识别的信号提取相关系统误差源，其定量评估有待后续工作完成。

\subsection{信号质量模型}
\label{subsec:syst-signal-model}

% 待填：改变 $J/\psi$ 信号模型——将 DSCB 替换为简单 Crystal Ball 或双高斯函数；将尾部参数 $\alpha_L$、$\alpha_R$ 在其拟合误差范围内变化；将 $\phi$ 的 Voigt 模型替换为相对论 Breit--Wigner 与 Crystal Ball 分辨率函数的卷积。
该误差源尚未完成定量评估，拟通过变动信号参数化模型来评估其对 $N_{sss}$ 的影响。

\subsection{本底质量模型}
\label{subsec:syst-background-model}

% 待填：改变本底模型——将 $J/\psi$ 的指数本底替换为一阶或二阶多项式；将 $\phi$ 的阈值幂律本底替换为指数函数或 Chebychev 多项式。
该误差源尚未完成定量评估，拟通过变动本底参数化形式来评估其对 $N_{sss}$ 的影响。

\subsection{固定拟合参数}
\label{subsec:syst-fixed-parameters}

% 待填：将固定的尾部幂次 $n_L$、$n_R$ 改变 $\pm 1$ 单位；将固定的 $\phi$ 宽度 $\Gamma_\phi$ 在其 PDG 误差范围内变化。
该误差源尚未完成定量评估，拟通过变动拟合中固定参数的值来评估其对 $N_{sss}$ 的影响。

\subsection{拟合范围}
\label{subsec:syst-fit-range}

% 待填：将 $J/\psi$ 和 $\phi$ 的质量拟合范围改变 $\pm 10\%$；在扩展和收缩范围内重新拟合，记录 $N_{sss}$ 的变化。
该误差源尚未完成定量评估，拟通过变动质量拟合窗口的范围来评估其对 $N_{sss}$ 的影响。

\subsection{拟合稳定性}
\label{subsec:syst-fit-stability}

% 待填：检查不同初始参数值下的拟合收敛性；利用伪实验进行系综检验，验证拟合偏差相对于统计误差可忽略。
该误差源尚未完成定量评估，拟通过伪实验系综检验来验证拟合程序的无偏性。

\subsection{sPlot 相关误差}
\label{subsec:syst-splot}

% 待填：评估由于拟合变量之间的相关性导致的 sWeight 计算中的潜在偏差；将 sWeight 加权的分布与更紧质量窗口选择得到的结果进行比较作为交叉检验。
该误差源尚未完成定量评估，拟通过与紧质量窗口选择结果的交叉检验来评估 sWeight 方法的潜在偏差。

\section{接受度和效率修正相关的系统误差}
\label{sec:syst-efficiency}

\subsection{蒙特卡罗统计误差}
\label{subsec:syst-mc-stat}

用于计算接受度和效率图的 MC 样本的有限规模在每个 $(p_{\mathrm T}, |y|)$ 分箱中引入了统计误差。该误差通过在 Clopper--Pearson 二项式误差范围内浮动逐分箱修正因子传播至修正后产额。

\subsection{接受度图分箱}
\label{subsec:syst-acceptance-binning}

% 待填：改变接受度图的分箱方案，记录对修正后产额的影响。
该误差源尚未完成定量评估，拟通过改变接受度图的分箱方案来评估其对修正后产额的影响。

\subsection{效率图分箱}
\label{subsec:syst-efficiency-binning}

% 待填：改变效率图的分箱方案，记录对修正后产额的影响。
该误差源尚未完成定量评估，拟通过改变效率图的分箱方案来评估其对修正后产额的影响。

\subsection{缪子重建与鉴别效率}
\label{subsec:syst-muon-eff}

% 待填：缪子重建和鉴别效率测量的系统误差——若可能，将基于 MC 的效率与数据中的 tag-and-probe 测量结果进行比较；将单缪子效率修正因子在其误差范围内变化。
该误差源尚未完成定量评估，拟通过数据驱动的 tag-and-probe 方法与 MC 效率的比较来评估缪子效率的系统误差。

\subsection{带电强子重建与鉴别效率}
\label{subsec:syst-kaon-eff}

% 待填：带电强子重建和鉴别效率的系统误差——利用控制样本比较 MC 与数据中的径迹重建效率。
该误差源尚未完成定量评估，拟通过控制样本比较 MC 与数据中的径迹重建效率来评估带电强子效率的系统误差。

\subsection{触发效率}
\label{subsec:syst-trigger-eff}

% 待填：HLT 触发效率的系统误差——改变用于推导触发效率的 MC 样本中的 SPS 与 DPS 混合比例；若可能，利用 tag-and-probe 方法进行数据驱动的触发效率测量作为比较。
该误差源尚未完成定量评估，拟通过改变 MC 样本中 SPS/DPS 混合比例及数据驱动方法来评估触发效率的系统误差。

\subsection{顶点拟合和事例级选择效率}
\label{subsec:syst-vertex-event-eff}

% 待填：四缪子顶点和六径迹初级顶点关联效率的系统误差——在信号压低控制区中将 MC 推导的效率与数据进行比较。
该误差源尚未完成定量评估，拟通过控制区数据/MC 比较来评估顶点拟合和事例级选择效率的系统误差。

\section{候选重建与选择相关的系统误差}
\label{sec:syst-candidate-selection}

\subsection{多重候选处理}
\label{subsec:syst-multiple-candidates}

% 待填：将基于 $p_{\mathrm T}$-求和的标称最佳候选选择与基于顶点拟合 $\chi^2$ 的替代规则进行比较；记录修正后产额的变化。
该误差源尚未完成定量评估，拟通过比较不同最佳候选选择规则来评估其对修正后产额的影响。

\subsection{最佳候选选择}
\label{subsec:syst-best-candidate}

% 待填：改变最佳候选判据，评估对事例产额和运动学分布的影响。
该误差源尚未完成定量评估，拟通过改变最佳候选判据来评估其对事例产额和运动学分布的影响。

\subsection{闭合检验残余偏差}
\label{subsec:syst-closure-bias}

% 待填：MC 闭合检验中观察到的任何残余非闭合均作为系统误差计入。
MC 闭合检验指在 MC 样本中，使用与数据相同的事例筛选、重建和产额修正方法，并将修正后的产额与已知的 MC 生成产额进行比较。任何残余的非闭合均作为系统误差计入。
该误差源尚未完成定量评估。

\section{样本建模相关的系统误差}
\label{sec:syst-sample-modeling}

\subsection{信号蒙特卡罗建模}
\label{subsec:syst-signal-mc-modeling}

% 待填：比较从不同 MC 产生子（如 \textsc{Pythia8} 与 \texttt{Helac-Onia}）推导的效率图；差异作为建模误差计入。
该误差源尚未完成定量评估，拟通过比较不同 MC 产生子推导的效率来评估信号建模的系统误差。

\subsection{运动学重加权}
\label{subsec:syst-kinematic-reweighting}

% 待填：对 MC 样本进行重加权以匹配数据的运动学分布，重新评估效率；与标称效率的差异作为系统误差计入。
该误差源尚未完成定量评估，拟通过对 MC 进行数据驱动的运动学重加权来评估其对效率修正的影响。

\subsection{不同子过程样本导致的效率修正变化}
\label{subsec:syst-model-dependence}

为评估效率修正对输入产生运动学的敏感性，本分析比较了由不同子过程MC样本分别构造的效率图。这里变化的是用于计算效率的子过程样本，而不是对同一物理模型参数进行连续变动。本分析考虑了五种效率情景：
\begin{itemize}
  \item \DPSone（基准）：两个$J/\psi+g$子过程组合形成的DPS样本；
  \item \DPStwoLO：LO色单态双$J/\psi$子过程与独立$gg\to gg$子过程组合形成的DPS样本；
  \item \DPStwoNLO：含额外实辐射胶子的NLO*色单态双$J/\psi$子过程与独立$gg\to gg$子过程组合形成的DPS样本；
  \item \SPSLO：LO色单态双$J/\psi$子过程的SPS样本；
  \item \SPSNLO：含额外实辐射胶子的NLO*色单态双$J/\psi$子过程的SPS样本。
\end{itemize}
每种情景均生成独立的接受度和效率图，并按照相同的程序计算效率修正后的信号产额，结果汇总于表~\ref{tab:eff-model-variations}。

\begin{table}[htbp]
  \centering
  \begin{threeparttable}[c]
    \caption{五种效率变动情景下的加权拟合产额。各情景均以独立的接受度和效率图进行修正后重新拟合。\SPSLO 情景因 SPS 截面理论预期极低而被排除（详见正文）。}
    \label{tab:eff-model-variations}
    \small
    \begin{tabular}{p{0.60\linewidth}cc}
      \toprule
      效率情景 & $N_{sss}^{\mathrm{corr}}$ & 相对变化（\%） \\
      \midrule
      \DPSone（基准） & $(9.35 \pm 1.94)\times10^4$ & --- \\
      \DPStwoLO & $(7.07 \pm 41.83)\times10^4$ & $-$24.4 \\
      \DPStwoNLO & $(5.83 \pm 1.07)\times10^4$ & $-$37.6 \\
      \SPSLO & $(24.67 \pm 5.57)\times10^4$ & +163.9 \\
      \SPSNLO & $(8.34 \pm 1.73)\times10^4$ & $-$10.8 \\
      \bottomrule
    \end{tabular}
  \end{threeparttable}
\end{table}

\SPSLO 情景给出的加权产额（$2.4672\times10^5$）远大于其他四种情景的结果，推测是因为该样本主要集中在低横动量区域，而使得高横动量区间的统计量相对较少，估计值产生较大偏差。相关多夸克偶素计算表明，纯SPS贡献相对于MPS贡献受到显著压低~\cite{Lansberg:2014swa}，因此该情景不作为标称包络的组成部分。以其余四种情景（\DPStwoLO、\DPStwoNLO、\SPSNLO 与基准\DPSone）相对于基准值的包络半宽定义当前已评定的效率修正系统误差分量：

\begin{equation}
\label{eq:model-uncertainty}
\Delta_{N}^{\mathrm{subproc}} =
\frac{9.3500\times10^4-5.8348\times10^4}{2}
=1.7576\times10^4
\quad (\text{相对误差 }18.8\%).
\end{equation}

该值被传播至修正后的信号产额和截面计算中，并在结果中标记为$\mathrm{syst,\,partial}$。该分量在所有运动学分箱中视为完全相关；由于其他系统误差源尚未全部完成定量评定，不能据此判断它是总系统误差中的主导来源。

\subsection{事例堆积建模}
\label{subsec:syst-pileup}

% 待填：对 MC 事例堆积分布进行重加权以匹配数据事例堆积分布；比较重加权前后的效率。
该误差源尚未完成定量评估，拟通过对 MC 事例堆积分布进行数据驱动的重加权来评估其对效率的影响。

\section{亮度与分支比误差}
\label{sec:syst-normalization}

\subsection{积分亮度}
\label{subsec:syst-luminosity}

数据样本的积分亮度为 $\SI{289.21 \pm 5.97}{\per\femto\barn}$，对应的相对误差为 2.1\%。该数值由本分析使用的 Muon 数据质量认证 JSON 文件通过 CMS 官方 \texttt{brilcalc} 工具计算得到；亮度标定和公开亮度结果参见 CMS 亮度公开页面~\cite{CMS:LUMI-PUB}。该误差在所有分箱中完全相关，并作为全局归一化误差处理。

\subsection{分支比}
\label{subsec:syst-branching-fractions}

截面计算采用第~\ref{sec:theory-quarkonium-production}~节列出的 $J/\psi \to \mu^+\mu^-$ 与 $\phi \to K^+K^-$ 世界平均分支比及其不确定度~\cite{PDG2024}。这些不确定度作为完全相关的归一化误差传播。

\section{系统误差的综合}
\label{sec:syst-combination}

\subsection{相关性假设}
\label{subsec:syst-correlations}

以下系统误差源在所有运动学分箱之间视为完全相关，因为它们代表全局归一化效应：积分亮度、分支比，以及不同子过程MC效率图导致的修正变化。信号提取相关的系统误差（信号和本底模型变动、拟合范围等）原则上可能随运动学分箱变化，但在此阶段先保守地视为分箱间完全相关。MC统计误差和分箱效应视为分箱间不相关。

\subsection{综合方法}
\label{subsec:syst-combination-procedure}

总系统误差由所有独立误差源平方求和开方得到，假定它们互不相关。对于同时包含相关和不相关分量的误差源，两部分先分别求和再合并。

\subsection{系统误差汇总表}
\label{subsec:syst-summary-table}

表~\ref{tab:syst-summary}汇总了当前已评估的所有系统误差源及其对修正后信号产额$N_{sss}^{\mathrm{corr}}$的相对贡献。标记为「待评估」的误差源尚未完成定量研究，将在后续工作中补充。

\begin{table}[htbp]
  \centering
  \begin{threeparttable}[c]
    \caption{各系统误差源对$N_{sss}^{\mathrm{corr}}$的相对贡献汇总。}
    \label{tab:syst-summary}
    \begin{tabular}{lcc}
      \toprule
      误差源 & 相对贡献（\%） & 相关性 \\
      \midrule
      \multicolumn{3}{l}{\textbf{已评估}} \\
      \quad 积分亮度~\cite{CMS:LUMI-PUB} & 2.1 & 完全相关 \\
      \quad $\mathcal{B}(J/\psi \to \mu^+\mu^-)$~\cite{PDG2024} & 0.55 & 完全相关 \\
      \quad $\mathcal{B}(\phi \to K^+K^-)$~\cite{PDG2024} & 1.0 & 完全相关 \\
      \quad 不同子过程MC效率图导致的修正变化（不含\SPSLO） & 18.8 & 完全相关 \\
      \midrule
      \multicolumn{3}{l}{\textbf{待评估}\tnote{*}} \\
      \quad 信号质量模型 & * & 待定 \\
      \quad 本底质量模型 & * & 待定 \\
      \quad 固定拟合参数 & * & 待定 \\
      \quad 拟合范围 & * & 待定 \\
      \quad 拟合稳定性 & * & 待定 \\
      \quad sPlot相关误差 & * & 待定 \\
      \quad MC统计误差 & * & 不相关 \\
      \quad 接受度图分箱 & * & 待定 \\
      \quad 效率图分箱 & * & 待定 \\
      \quad 缪子重建与鉴别效率 & * & 待定 \\
      \quad 带电强子重建与鉴别效率 & * & 待定 \\
      \quad 触发效率 & * & 待定 \\
      \quad 顶点拟合与事例级选择效率 & * & 待定 \\
      \quad 多重候选处理 & * & 待定 \\
      \quad 最佳候选选择 & * & 待定 \\
      \quad 闭合检验残余偏差 & * & 待定 \\
      \quad 信号MC建模 & * & 待定 \\
      \quad 运动学重加权 & * & 待定 \\
      \quad 事例堆积建模 & * & 待定 \\
      \bottomrule
    \end{tabular}
    \begin{tablenotes}
      \item[*] 尚未完成定量评估的误差源。
    \end{tablenotes}
  \end{threeparttable}
\end{table}
